% YZM 5110 Multimodal Learning, Project 1 -- research report (IEEE conference format)
% Build: pdflatex report (twice, so the citation numbers resolve)
% Needs IEEEtran.cls (TeX Live package texlive-publishers). The bibliography is embedded: no bibtex run needed.
% report_extended.tex is the long version of this report with every intermediate table and paragraph.
\documentclass[conference]{IEEEtran}

\usepackage[T1]{fontenc}
\usepackage[utf8]{inputenc}
\usepackage{lmodern}
\usepackage{microtype}
\usepackage{booktabs}
\usepackage{tabularx}
\usepackage{longtable}
\usepackage{array}
\usepackage{amsmath,amssymb}
\usepackage{url}
\usepackage[hidelinks]{hyperref}

\sloppy
\emergencystretch=2em
\newcommand{\code}[1]{\texttt{\detokenize{#1}}}

\begin{document}

\title{Does MAMEX's Modality Gate Trust Corrupted Text?\\
A Stress Test and a Visual-Grounding Fix}

\author{\IEEEauthorblockN{Hatice Sebla Karabunar}
\IEEEauthorblockA{YZM 5110 Multimodal Learning, Project 1 (Dr.\ Ekrem \c{C}etinkaya), November 2026\\
Repository: \url{https://github.com/karabunar/mamex-robustness-analysis}\\
Commit / tag: \texttt{v1.0} (to be set at submission)\\
AI use: disclosed in Appendix~\ref{app:ai}}}

\maketitle

\begin{abstract}
MAMEX is a cold-start recommender whose gate weights an image expert and a text expert. We ask what the gate does when an item's text is present but misleading. On Amazon Baby we reproduce the released code (Recall@20 0.0776 against 0.0771 published; Recall@10 stays about 0.0025 short and the cause is not identified). Corrupting the text of 40\% of the cold items at test time lowers Recall@20 by 14\% and the gate keeps trusting the text. A parameter-free fix, Visual Grounding Regularization (VGR), recovers 27--38\% of the loss, but about half of that is also reached by controls that ignore the image. A training-free repair that blends in the text of the visually nearest warm items does better (M3: +0.0088 Recall@20 over ten seeds, about 85\% of the loss on five seeds), beats random-neighbour and permuted-score controls, and keeps its ranking on a second dataset, where it recovers less (about 38\%). The gate's preference for text on cold items follows from an item having no training interactions, not from being cold. The distributed features are LATTICE's, not CLIP.
\end{abstract}

\begin{IEEEkeywords}
cold-start recommendation, multimodal recommendation, mixture of experts, robustness, reproducibility
\end{IEEEkeywords}

%====================================================================
\section{Research Question, Scope, Contributions}

\textbf{Question.} MAMEX~\cite{nguyen2025mamex} builds cold-item embeddings by letting a gate weight an image expert and a text expert. When an item's text is present but misleading (clickbait-like), does the gate keep trusting it, and can a parameter-free image--text agreement signal reduce the damage without hurting clean accuracy?

\textbf{Scope.} Amazon Baby item cold start (7,050 items, 1,410 cold, 19,442 training users). Text is corrupted at test time only, at the embedding level, on a share of the cold items. Three seeds for the first experiments, ten for the main repair comparison. Amazon Sports, built by us from the LATTICE files (not the authors' split), is a robustness check only. Not in scope: real clickbait text, other cold-start models.

\textbf{Motivation.} None of the papers surveyed in Section~\ref{sec:lit} reports performance when a modality is present but misleading.

\textbf{Contributions.} (1)~a measured reproduction of the MAMEX Baby baseline, with an investigation of the remaining gap; (2)~a stress-test protocol for corrupted cold-item text; (3)~VGR and a training-free neighbour repair, each with controls that isolate where the effect comes from; (4)~measurements of why the gate prefers text on cold items; (5)~a tested, packaged codebase with exact commands.

Every number is measured in the experiments of Section~\ref{sec:design} or taken from a cited paper; what a paper reports is attributed to it, an inference is marked as our reading. Every hypothesis, with its prediction written before the run and its outcome, is in \code{docs/hypotheses.md} of the repository.

%====================================================================
\section{Literature Search and Critical Synthesis}
\label{sec:lit}

\textbf{Procedure.} Web search on 2026-10-16: one query per paper by title (VBPR, MMGCN, CLCRec, MoToRec, MAMEX) and two topic queries (multimodal cold-start with adaptive mixture of experts; multimodal tokenization cold-start). DGMRec surfaced as a follow-up. MILK~\cite{bai2024milk} was added because the acknowledgements of the MAMEX README name it and it is the strongest baseline in MAMEX's own table. This is a targeted selection, not a systematic review. Selection criteria: the paper that introduced the base idea (VBPR, MMGCN), the paper reproduced (MAMEX), competing cold-start philosophies (CLCRec: contrastive alignment; MoToRec: discrete tokens), and work on a related failure mode (DGMRec, MILK: missing modalities). Five papers form the matrix, DGMRec and MILK are supporting papers. All were read in full; each matrix row was checked against the text (MILK's tables, which text extraction garbles, against the claims in its text).

\begin{table*}[!t]
\centering\footnotesize
\caption{Condensed comparison; results as reported by each paper. Cold-start definitions differ, so numbers are not comparable across rows. The course's seven-field matrix is Appendix~\ref{app:matrix}.}
\label{tab:lit}
\renewcommand{\arraystretch}{1.1}
\begin{tabularx}{\textwidth}{@{}p{1.8cm}p{3.0cm}XXp{2.6cm}@{}}
\toprule
Paper & Cold-start definition & Modalities and fusion & Strongest reported evidence & Misleading input tested? \\
\midrule
VBPR~\cite{he2016vbpr} & items with $<5$ training positives & image (4096-d CNN) as extra factors in BPR-MF & Women, cold-start AUC 0.6813 vs 0.5281 for BPR-MF & no (image only) \\
MMGCN~\cite{wei2019mmgcn} & none: items must be in the graph & visual, acoustic, text; one GCN per modality, summed & Recall +8.76 to +15.59\% over the best baseline & no (calls text ``noisy, incomplete, and even irrelevant'', does not test it) \\
CLCRec~\cite{wei2021clcrec} & items unseen in training & visual, acoustic, text; contrastive objectives & gains over DUIF, DropoutNet, MTPR, CB2CF, Heater; single runs & no \\
MoToRec~\cite{liu2026motorec} & items with $<10$ training interactions & BEiT image, BGE text; discrete tokens (RQ-VAE), LightGCN & Baby cold-start R@20 0.0347 & no (argues noisy embeddings hurt, does not corrupt inputs) \\
MAMEX~\cite{nguyen2025mamex} & complete item cold start & CLIP image and text (as stated); linear-expert MoE and gate & Baby Rec@20 0.0771, NDCG@20 0.0363 & no (clean features only) \\
DGMRec~\cite{kim2025dgmrec} & missing modalities; new items & CNN image, SBERT text; disentangle, then generate the missing feature & at 80\% missing, drops 10.2\% vs 15.1\% for the baseline compared & missing: yes; misleading: no \\
MILK~\cite{bai2024milk} & new items with one modality missing at test & visual and text; Dirichlet-weighted fusion environments, variance penalty & Baby, one modality missing: Rec@20 0.0628, NDCG@20 +11.2\% over the best baseline & missing (mean-filled): yes; misleading: no \\
\bottomrule
\end{tabularx}
\end{table*}

\textbf{Synthesis (our reading).}
\begin{enumerate}
\item \emph{Shared pipeline, incomparable numbers.} All papers use frozen pre-extracted content features and learn what sits on top; content helps items with little history. Cold means fewer than 5 positives (VBPR), fewer than 10 (MoToRec) or none (CLCRec, MAMEX), and MMGCN and MoToRec split interactions at random, which cannot evaluate cold items. MoToRec's Baby R@20 0.0347 (cold group) and MAMEX's 0.0771 come from different protocols.
\item \emph{The nearest remedies treat a failing modality as absent.} MILK trains with Dirichlet-drawn fusion weights, cyclically shifted so that each modality dominates in some environment, and penalises the variance of the loss across them; it mean-fills a missing modality (zeros, means or a learned map did not help: Baby Rec@20 0.0501, 0.0525, 0.0507 against 0.0628). DGMRec generates the missing feature. Both are training-time remedies; VGR and our repair act on a trained model. The settings differ (absent against misleading text), so we draw no numerical comparison.
\item \emph{The gap.} None of the seven papers tests a modality that is present but misleading (the search was targeted, so other work may exist). That is the question we ask of MAMEX's gate.
\end{enumerate}

%====================================================================
\section{The Method: MAMEX as Implemented}

We read the model code and checked it against the paper's equations.
\begin{itemize}
\item Frozen, L2-normalised image and text features each pass through a mixture-of-experts layer (4 linear experts, top-2 by a learned gate), giving a 128-d embedding per modality and a load-balancing loss.
\item The two embeddings are concatenated and a fusion gate (MLP, ReLU, dropout, two outputs, softmax) gives $[w_v, w_t]$; the item embedding is $w_v e_{\text{img}} + w_t e_{\text{txt}}$. There is no item-id embedding, so an item's embedding is a function of its features alone.
\item User embeddings are free parameters. The loss is BPR plus L2, an alignment term (MSE between item and modality embeddings), the expert balance loss and a KL term keeping the mean fusion gate near uniform. Cold items are scored by the inner product of the user embedding with an embedding built from their features only.
\item \emph{Against the paper} (arXiv 2508.08042v1, Eqs.~3--11): the structure matches. Two differences: Eq.~5 describes a sparse softmax gate while the code uses an MLP and an ordinary softmax over two outputs; the text after Eq.~8 calls the alignment loss an MSE but the equation is a squared L2 norm, and the code's MSE is smaller by the embedding size (128). The paper's weight $\lambda_1 \in \{1, 0.1, 0.01\}$ does not say which scale it uses.
\end{itemize}

%====================================================================
\section{Reproduction Target, Implementation, Deviations}
\label{sec:repro}

\textbf{Target} (published, Table~1 of the paper, kept apart from measured numbers): Rec@10 0.0512, Rec@20 0.0771, NDCG@10 0.0292, NDCG@20 0.0363. \textbf{Implementation.} The upstream code (commit \code{88a932e}) was restructured into a package with characterisation tests before and unit and end-to-end tests after (158 tests); training and evaluation logic are unchanged. Added: test-time corruption, VGR and the repairs with controls, $\gamma$ selection on validation data, the evaluation report, and an optional training-time corruption for the ablation (off by default). A reused-component inventory is in \code{docs/reused_components.md}.

\textbf{Deviation 1: the features are LATTICE's, not CLIP.} The paper and README describe CLIP features; the authors host only the processed data, and an open issue (\#2, August 2025) asks for the preprocessing; the reply (19 September 2025) points to MMSRec~\cite{mmsrec} and McAuley's raw data and names no extraction script. The distributed files are not CLIP-shaped: images are 4,096-d, non-negative and 79.2\% exactly zero (what a ReLU fully connected layer of an ImageNet CNN gives, as VBPR's fc7 features~\cite{he2016vbpr}); CLIP image embeddings have 512 to 1,024 dimensions and are dense and signed~\cite{radford2021clip}. Text is 1,024-d, dense and signed. We then compared the Baby folder published by the LATTICE authors~\cite{zhang2021lattice} (files dated September 2021): \code{image_feat.npy} and \code{text_feat.npy} are identical to the distributed ones (7,050 rows in the same order, maximum absolute difference 0.0). The data folder of our repository is byte-identical (SHA-256, 30 of 30 files) to the authors' \code{processed\_data\_baby.zip}, and the Baby interactions equal the union of LATTICE's 5-core train, validation and test pairs (160,792 pairs). So every result here uses LATTICE's features. We do not know why the MAMEX paper says CLIP, and we do not offer this as the cause of the Rec@10 gap (Rec@20 and NDCG@20 match, and no experiment varied the features). The reproduction target is therefore ``the released code on the released data''.

\textbf{Other deviations.} Batch size, epochs, embedding size and hardware are unstated in the paper; repository defaults are used (batch 2,048, 128-d, up to 1,000 epochs, early stop 200). The best epoch is chosen on clean validation NDCG@10. Seeds 2025--2027 for the first experiments (2025 is the upstream default, the others are consecutive integers; none was chosen by looking at results). The reproduction criterion (the claim is questioned for a metric if the published value lies outside the seed range) was written after the baseline was seen.

%====================================================================
\section{Proposed Improvement and Rationale}

\textbf{VGR.} After the gate, $w_t' = w_t \cdot \sigma(\gamma \cos(e_{\text{img}}, e_{\text{txt}}))$, then $[w_v, w_t']$ is renormalised. The cosine uses the 128-d post-MoE embeddings and carries no gradient; there are no learned parameters, so VGR is switched on or off on the same clean-trained checkpoint. \emph{Rationale} (stated before results): alignment pulls the two embeddings together on clean data, so a low cosine signals text that disagrees with the image. A cost known in advance: $\sigma<1$ also down-weights clean text.

\textbf{Training-free repair.} VGR lowers the weight of a suspect text but puts nothing in its place, so we tested replacing part of it by $t_{\mathrm{knn}}$, the mean unit text of the ten nearest warm items by image cosine. \emph{M1} blends it into every cold item, $t'=(1-\lambda)t+\lambda t_{\mathrm{knn}}$. \emph{M3} blends it only where VGR's cosine $c_i$ is low: $\lambda_i=\lambda_{\max}\sigma(g(Q_q-c_i))$, with $Q_q$ the $q$-quantile of the cosines over the cold items. (An earlier variant, M2, gated by the agreement between $t$ and $t_{\mathrm{knn}}$, was dropped: its per-item score adds nothing over a permuted score.) Controls: random warm neighbours, and permuted scores. All parameters ($\gamma$, $\lambda$, $\lambda_{\max}$, $q$, $g$) are chosen on noisy validation data, never on the test set.

\textbf{Hypotheses.} Each was written, with the result that would falsify it, before its runs (H1--H4 before the first full results; the repository history was built after the experiments, so the order is our statement, not a commit timestamp). Table~\ref{tab:hyp} lists all of them with their outcomes.

\begin{table*}[!t]
\centering\footnotesize
\caption{Hypotheses and outcomes. ``Supported'' means all stated predictions held; details and the exact thresholds are in \code{docs/hypotheses.md}.}
\label{tab:hyp}
\renewcommand{\arraystretch}{1.05}
\begin{tabularx}{\textwidth}{@{}l>{\raggedright\arraybackslash}p{6.7cm}>{\raggedright\arraybackslash}X@{}}
\toprule
 & Prediction & Outcome \\
\midrule
H1 & corrupted text lowers accuracy and the gate does not move away from text & supported: Recall@20 $-14\%$; $w_t$ 0.550 $\to$ 0.563 \\
H2 & VGR recovers part of the loss & partly: 27--38\% of the loss, 3/3 seeds \\
H3 & VGR costs under one seed std of clean accuracy and under 5\% latency & quality yes; latency 5.6\%, limit missed \\
H4 & the gain needs the image (controls recover less) & partly: controls recover about half \\
H5, H25 & noise-augmented training raises noisy accuracy (rate 0.4; then rates 0.1, 0.2) & not supported at 0.4 ($-0.0100$ Recall@20); at 0.1 it helps ($+0.0028$, $+0.0070$), at 0.2 it does not \\
H6, H7 & the collapse is on corrupted warm items; cold items lie far from the warm centroid & not supported: it is on cold items; they lie closer (ratio 0.72) \\
H8, H9 & image-guided repairs beat controls; M3 (flagged items) beats M1 & M1/M3 beat controls; M2's score adds nothing; M3 above M1 by little \\
H10, H11 & M3 holds under other noise, mismatched validation and five seeds & supported, except M3 $\ge$ M1 in only 5 of 9 settings (M1 better at rate 0.6) \\
H12, H14 & the Rec@10 gap comes from the epoch rule, candidate set or per-metric epochs & no validation rule, candidate set or per-metric maxima closes it; a test-chosen epoch closes about 60\% \\
H13 & gate bias or the norm of cold expert outputs drives the collapse & not supported; the image-side outputs of cold items do \\
H15 & at least 70\% of M3's gain is on corrupted items & mean 72\%; below 70\% in two of five seeds \\
H16 & M3 beats M1 over ten seeds, mean difference $+0.0005$ to $+0.0025$ & supported: $+0.0019$, interval $[+0.0004,+0.0037]$ \\
H17 & items without training interactions behave like cold items & supported: $w_t$ 0.96 against 0.35 on trained items \\
H18 & the method ranking holds on a second dataset & supported in all five predictions; effect size smaller \\
H19, H20 & $w_t$ falls with interaction count; the user direction of the image output carries it & H19 supported (control just missed); H20: user direction carries information beyond the count, not the strongest correlate \\
H24 & a $\gamma$ rule that maximises validation VGR minus shuffled control picks $\gamma\ge2$ and raises the cosine-specific gain & partly: picks 4, 1, 2; cosine-specific gain up in 2 of 3 seeds, total gain down; released rule kept \\
H26 & a 20-dimensional user subspace carries most of what the gate reads & no: about half (50\%, 35\%) \\
H21 & swapping the user-direction component of the image output moves $w_t$ by at least 0.2, random directions by under 0.05 & random: yes; cold $-0.38$ and $-0.18$ (second seed below 0.2); warm $+0.20$, $+0.20$ \\
H22 & the neighbour text is less informative on Sports than on Baby & keeps 76\% of clean Recall@20 against 93\% (cosine differs only 0.02) \\
H23 & M3's parts add up over random partitions & no: the interaction is as large for random partitions as for corrupted/clean \\
\bottomrule
\end{tabularx}
\end{table*}

%====================================================================
\section{Experimental Design}
\label{sec:design}

\begin{itemize}
\item \textbf{Data.} Baby item cold-start split as distributed: cold items never appear in training; validation and test contain cold items only (705 items each); candidates are all 1,410 cold items.
\item \textbf{Noise protocol} (a design decision). Test time only, on the cold items' text features; training and validation stay clean, so the clean control is exact and the gate cannot learn to hide the weakness. \emph{Hashtag} noise (main): a corrupted item's unit text vector is mixed (weight 0.8) with one of 5 shared ``clickbait'' prototypes drawn from warm items. Also: \emph{shuffle} (another item's text) and \emph{gaussian}. Rate 0.4 (564 of 1,410 items). The noise is simulated on embeddings; it is not real clickbait text. The same seed gives the same mask in every mode.
\item \textbf{Conditions.} Clean; stress test (noisy); VGR ($\gamma$ from $\{1,2,4,8,16\}$ chosen on noisy validation data with a different mask); clean text with VGR; shuffled-cosine and constant-factor controls; the repairs.
\item \textbf{Metrics.} Recall and NDCG at $K\in\{10,\dots,50\}$, headline Recall@20 and NDCG@20; mean $\pm$ sample std and per-seed paired differences. Success@$K$ and MRR with bootstrap intervals over test users. No significance test with $n=3$.
\item \textbf{No tuning on the test set.} $\gamma$ and the epoch use validation data only; a test-set $\gamma$ sweep is a sensitivity table reported after selection.
\item \textbf{Environment.} One RTX A4000; about 19 min per training run alone. PyTorch 2.9.1 is pinned and reproduces the saved baselines bit for bit; the system PyTorch 2.14.1, installed later, does not (Appendix~\ref{app:record}). The core results were made before that change; the ablation and probe runs (H5--H7) used 2.14.1 with a clean-trained control from the same environment, and H11 onward ran in 2.9.1.
\end{itemize}

%====================================================================
\section{Results, Controls, Uncertainty, Failure Analysis}
\label{sec:results}

Test split, cold items, hashtag noise at 0.4, mean over seeds 2025--2027 unless stated.

\subsection{Reproduction}

\begin{table}[!t]
\centering\footnotesize
\caption{Published vs measured clean baseline.}
\label{tab:repro}
\begin{tabular}{@{}lcccc@{}}
\toprule
 & Rec@10 & Rec@20 & NDCG@10 & NDCG@20 \\
\midrule
Published & 0.0512 & 0.0771 & 0.0292 & 0.0363 \\
Measured, 3 seeds & 0.0487 & 0.0776 & 0.0285 & 0.0365 \\
Std over seeds & 0.0023 & 0.0042 & 0.0017 & 0.0022 \\
\bottomrule
\end{tabular}
\end{table}

Rec@20, NDCG@10 and NDCG@20 lie inside the seed range and within 0.001 of the mean; Rec@10 is 0.0025 above the mean. Four diagnostics, each with its rule written first, did not close the gap. \emph{D1 seeds:} ten seeds give Rec@10 $0.0488\pm0.0004$ (standard error), the published 0.0512 lies outside two standard errors and none of ten seeds reaches it (best 0.0507). \emph{D2 alignment-loss scale} (weights 1.28, 12.8, 128, the paper's three values): none brings Rec@10 within 0.001 of the published value (closest 0.0500). \emph{D3 epoch rule (H12):} choosing the epoch by validation Recall@10, NDCG@20 or Recall@20 gives 0.0491, 0.0487 and 0.0482; choosing it by test Recall@10 (a diagnostic only) gives 0.0502, about 60\% of the gap but still 0.0010 short. \emph{H14:} ranking among the 705 test items only raises every metric by 53--59\% (Recall@20 0.0776 to 0.1190), which excludes this protocol, and the per-metric best epochs match the published vector in no seed. So the gap is \emph{not identified}. The one candidate left is an epoch chosen with the test set in the published run, which we cannot check.

\subsection{Stress test and VGR}

\begin{table}[!t]
\centering\footnotesize
\caption{Conditions, hashtag noise at 0.4 (mean $\pm$ std, 3 seeds).}
\label{tab:cond}
\begin{tabular}{@{}lcc@{}}
\toprule
Condition & Recall@20 & NDCG@20 \\
\midrule
Clean baseline & $0.078\pm0.004$ & $0.037\pm0.002$ \\
Stress test (noisy text) & $0.067\pm0.002$ & $0.031\pm0.001$ \\
VGR (noisy text) & $0.070\pm0.002$ & $0.032\pm0.001$ \\
VGR (clean text) & $0.077\pm0.005$ & $0.036\pm0.002$ \\
Control: shuffled cosine (noisy) & $0.068\pm0.002$ & $0.031\pm0.001$ \\
Control: constant factor (noisy) & $0.068\pm0.002$ & $0.031\pm0.001$ \\
\bottomrule
\end{tabular}
\end{table}

\textbf{H1 holds.} Corruption lowers Recall@20 from 0.0776 to 0.0668 ($-14\%$) and NDCG@20 from 0.0365 to 0.0307. On the corrupted items $w_t$ goes from 0.550 to 0.563: the gate does not move away from text. \textbf{H2 holds, small.} Noisy VGR beats noisy in 3/3 seeds ($+0.0033$ Recall@20, 27--38\% of the loss) and moves $w_t$ to 0.416. \textbf{H3:} clean VGR changes Recall@20 by $-0.0004$ (below the seed std), but the forward pass slows by 0.36 ms, 5.6\%, above the 5\% limit set in advance. \textbf{H4:} at the selected $\gamma=1$ the shuffled-cosine and constant-factor controls recover about half of VGR's gain ($+0.0017$, $+0.0015$ against $+0.0033$); VGR beats both in 3/3 seeds, so the grounding-specific part is about $+0.0017$, about 15\% of the loss.

\emph{Uncertainty.} Paired bootstrap over the 10,474 test users (1,000 draws; user sampling only, not seed variance): noisy minus clean $-0.0109$ [$-0.0129,-0.0087$]; VGR minus noisy $+0.0033$ [$+0.0024,+0.0044$]; VGR minus shuffled $+0.0017$ [$+0.0011,+0.0023$]; clean VGR minus clean $-0.0004$ [$-0.0009,+0.0002$]. Success@$K$ and MRR agree with Recall and NDCG; VGR's gain appears from Success@5 upward, not at Success@1. \emph{Failure analysis} (seed 2025): users who see a corrupted test item gain $+0.0101$ Recall@20 from VGR (238 helped, 117 hurt), others lose $-0.0016$ (81 helped, 187 hurt); seeds 2026 and 2027 agree. \emph{Sensitivity} (seed 2025, post hoc): the selected $\gamma$ was the grid minimum. The total gain rises as $\gamma$ falls ($+0.0011$ at 16, $+0.0054$ at 0.25), but the cosine-specific part peaks at $\gamma=2$ and shrinks to $+0.0007$ at 0.25, where VGR is a uniform muting of text. VGR helps at every noise rate from 0.2 to 0.8 ($+0.0018$ to $+0.0070$). We did not choose $\gamma$ from this table, which would be tuning on the test set.

\subsection{Training-free repair}

\begin{table}[t]
\centering\caption{Recall@20, hashtag 0.4, seeds 2025--2027; gain over the unrepaired model (seeds with a positive gain).}
\label{tab:repair}
\small\setlength{\tabcolsep}{3pt}
\begin{tabular}{@{}>{\raggedright\arraybackslash}p{3.3cm}ccc@{}}
\hline
Method & noisy text & gain & clean text \\
\hline
Unrepaired & 0.0668 & --- & 0.0776 \\
VGR & 0.0701 & $+0.0033$ (3/3) & 0.0773 \\
M1 blind blend & 0.0746 & $+0.0078$ (3/3) & 0.0769 \\
M1, random neighbours & 0.0666 & $-0.0002$ (1/3) & 0.0759 \\
M3, cosine-gated & 0.0761 & $+0.0094$ (3/3) & 0.0792 \\
M3, random neighbours & 0.0644 & $-0.0023$ (0/3) & 0.0721 \\
M3, permuted cosine & 0.0724 & $+0.0056$ (3/3) & 0.0794 \\
M3, then VGR & 0.0772 & $+0.0105$ (3/3) & 0.0789 \\
Oracle (true mask) & 0.0735 & $+0.0068$ (3/3) & 0.0776 \\
\hline
\end{tabular}
\end{table}

\textbf{H8.} The image guidance is what helps: both repairs beat their random-neighbour controls in 3/3 seeds. M2's per-item agreement score does not: it loses to its permuted control in all three seeds, and its AUROC for finding the 564 corrupted items is only 0.51--0.71, against 0.999 for VGR's cosine. The blind blend M1 is the best single method of that round ($+0.0078$) at a clean-text cost of only $-0.0007$, which says that the nearest warm items' text is almost as informative as an item's own text for these features. \textbf{H9.} M3 gains $+0.0094$ (about 86\% of the 0.0109 lost; with VGR $+0.0105$, about 96\%), above M1 in all three seeds but by little ($+0.0029$, $+0.0005$, $+0.0011$), above both controls, and it changes clean text by $+0.0016$. The ``oracle'' that repairs exactly the corrupted items is not an upper bound: M3 and even M1 beat it. Building the repair costs 0.6 s once, the cosine 0.017 s, 41 MB, and the forward pass +2.4\%.

\begin{table}[t]
\centering\caption{Recall@20 of the unrepaired model and gains of VGR, M1, M3 and M3 then VGR (mean over three seeds; every gain positive in all three seeds).}
\label{tab:grid}
\footnotesize\setlength{\tabcolsep}{3pt}
\begin{tabular}{@{}lccccc@{}}
\hline
Test noise & unrep. & VGR & M1 & M3 & M3+VGR \\
\hline
hashtag 0.2 & 0.0716 & $+0.0015$ & $+0.0049$ & $+0.0060$ & $+0.0061$ \\
hashtag 0.4 & 0.0668 & $+0.0033$ & $+0.0078$ & $+0.0094$ & $+0.0105$ \\
hashtag 0.6 & 0.0599 & $+0.0042$ & $+0.0135$ & $+0.0101$ & $+0.0129$ \\
shuffle 0.2 & 0.0681 & $+0.0028$ & $+0.0065$ & $+0.0094$ & $+0.0094$ \\
shuffle 0.4 & 0.0604 & $+0.0072$ & $+0.0128$ & $+0.0152$ & $+0.0159$ \\
shuffle 0.6 & 0.0485 & $+0.0101$ & $+0.0236$ & $+0.0224$ & $+0.0241$ \\
gaussian 0.2 & 0.0736 & $+0.0009$ & $+0.0039$ & $+0.0048$ & $+0.0050$ \\
gaussian 0.4 & 0.0721 & $+0.0009$ & $+0.0043$ & $+0.0039$ & $+0.0047$ \\
gaussian 0.6 & 0.0664 & $+0.0024$ & $+0.0077$ & $+0.0075$ & $+0.0079$ \\
\hline
\end{tabular}
\end{table}

\textbf{H10 (Table~\ref{tab:grid}).} M3 beats the unrepaired model in all nine settings and all three seeds, and its mean beats VGR's. It beats both controls in every setting, though for Gaussian 0.4 only by $+0.0004$ over the permuted-cosine control. It is at least as good as M1 in only five of nine settings: M1 is better at rate 0.6 for every noise type, so with much corrupted text flagging items adds nothing over repairing all of them. With validation noise fixed at hashtag 0.4 (four mismatch settings) M3 keeps 99--116\% of its matched gain, so the repair does not depend on knowing the noise type or rate. The AUROC of the cosine is at least 0.993 everywhere.

\begin{table}[t]
\centering\caption{Recall@20, hashtag 0.4: unrepaired score and mean gains. Baby: ten seeds (2025--2034); Sports: three seeds (our split).}
\label{tab:seeds}
\scriptsize\setlength{\tabcolsep}{2pt}
\begin{tabular}{@{}lcccccc@{}}
\hline
 & unrep. & VGR & M1 & M3 & M3+VGR & M3$-$M1 \\
\hline
Baby, 10 seeds & 0.0668 & $+.0030$ & $+.0069$ & $+.0088$ & $+.0095$ & $+.0019$ \\
Sports, 3 seeds & 0.0775 & $+.0018$ & $+.0023$ & $+.0068$ & $+.0065$ & $+.0046$ \\
\hline
\end{tabular}
\end{table}

\textbf{H11, H16: ten seeds (Table~\ref{tab:seeds}).} M3 beats the unrepaired model in 10 of 10 seeds and is at least M1 in 8 of 10 (exceptions 2028, $-0.0001$; 2030, $-0.0019$). M3 minus M1 has mean $+0.0019$ (std 0.0029), paired bootstrap interval over seeds $[+0.0004, +0.0037]$ (10,000 resamples): it excludes zero but its lower end is close to it. One seed carries the mean (without seed 2031, $+0.0089$, it is about $+0.0010$; the median is $+0.0011$), and in seed 2030 M3 does not beat both controls. We read M3 as a small and consistent advantage over the blind blend, not a large one. On the first five seeds M3 recovers about 85\% of the loss, M3 then VGR about 93\%. \textbf{H15.} The corrupted items alone give 87, 54, 77, 62 and 71\% of M3's gain (mean 72\%); the clean-only gain is positive in every seed ($+0.0007$ to $+0.0030$; seed 2028 is $+0.00301$ against our $+0.003$ limit), so part of the benefit is neighbour text on clean items. (H23: the two parts do not add up to the full gain, and random partitions of the cold items show interactions as large, up to 0.0025, so this is a general property of ranking and the shares carry an uncertainty of about 0.002.) \textbf{H18: second dataset} (Sports: 35,598 users, 18,357 items, 3,671 cold items; our construction in the format of the MAMEX split). All five predictions hold: corruption lowers Recall@20 (0.0775 unrepaired, loss 0.0179), and VGR, M1 and M3 gain in all three seeds; M3 beats both controls ($\ge+0.0050$), VGR and M1 in all three. The size does not carry over: M3 recovers about 38\% of the loss (about 85\% on Baby), M1 about 13\%, and VGR after M3 adds nothing. The ranking (M3 above M1 above VGR) is the same. H22: the neighbour text is less informative on Sports: replacing every cold item's text by it keeps 76\% of the clean Recall@20, against 93\% on Baby, although its cosine to the item's own text differs only by 0.02 (0.683, 0.707); this is consistent with the smaller recovery, but we did not show that it explains all of it.

\subsection{Why the gate prefers text on cold items}

Training on corrupted warm items (H5, rate 0.4 in every epoch) did not help: noisy Recall@20 is 0.0568, below the stress-test model's 0.0668 in 3/3 seeds, and VGR adds only $+0.0003$. The cause is the gate: on cold items it puts $w_t=0.917$ on noisy text (0.986 on clean) against 0.563 and 0.550 without augmentation. The failure belongs to the high rate (H25, two seeds each): at rate 0.2 the gate is nearly as collapsed ($w_t$ 0.83 on clean cold items) and noisy Recall@20 is no better than without augmentation (0.0629, 0.0652 against 0.0667, 0.0652); at rate 0.1 the gate is less collapsed (0.62, 0.68) and noisy Recall@20 is above the clean-trained model's in both seeds (0.0695, 0.0722; $+0.0028$, $+0.0070$), though still below the training-free repair M3 (0.0780, 0.0731). The collapse grows with the rate. We then tested where it comes from.
\begin{itemize}
\item \emph{Cold items only (H6, H7).} On warm items the augmented gate gives 0.44 ($w_t$ 0.31 on freshly corrupted ones) but 0.99 on cold items. Cold items do not lie farther from the warm centroid (ratio 0.72, not above 1.2), and distance does not correlate with $w_t$ (Spearman $+0.08$ to $+0.15$).
\item \emph{Not bias, not norm (H13).} With both expert outputs set to zero the text-minus-image logit is only $+0.18$ and $+0.22$ (we predicted above $+3$). Rescaling cold items' outputs to the warm norm leaves $w_t$ at 0.99. Replacing both outputs by a random warm item's brings it to the warm level (0.436); replacing only the image output gives 0.32 and 0.36, only the text output 0.95 and 0.91 (exploratory). The image-side outputs of cold items drive the text preference.
\item \emph{Missing training interactions (H17).} We removed the training interactions of 10\% of the warm items (564 items, kept in the item set). In the augmented model they get $w_t$ 0.966 and 0.963, like cold items (0.986, 0.990) and unlike trained items (0.353, 0.361); the clean-trained control stays between 0.48 and 0.54 in all groups. So the preference follows from an item not being in the training loss, not from being cold. (Our hypothesis text wrongly said cold items are absent from negative sampling; negatives are drawn over all items.)
\item \emph{A graded trace (H19, H20).} Among warm items $w_t$ falls with the log training interaction count (Spearman $-0.69$ and $-0.76$; control $-0.39$ and $-0.32$, just missing our below-0.30 prediction), but the least-interacted fifth (median 5 interactions) is at only 0.54--0.62, far below 0.96 for items with none, so the change near zero is abrupt. The image output's cosine to the mean user embedding correlates at $-0.69$ and $-0.61$ and keeps $-0.50$ and $-0.35$ after removing the count; the density in the raw feature space does not ($+0.05$, $+0.08$). The norm is the strongest raw correlate but H13 showed it is not what the gate uses. H21 is an intervention on this component: giving every cold item the warm mean of its user-direction component lowers $w_t$ by 0.38 and 0.18 (seeds 2025, 2026; the second below our 0.2), giving warm items the cold mean raises it by 0.20 and 0.20, and the same swap along random directions changes it by 0.001; in the clean-trained models it changes it by at most 0.05. So the gate does read the user direction, one dimension of 128, though not alone (cold items stay at 0.61 and 0.80). H26 widens the swap to the top 5 and top 20 user directions: the drop is only 0.28 and 0.15 (k=5), 0.34 and 0.22 (k=20), about half of the 0.62--0.68 that replacing the whole image output gives, and random subspaces change $w_t$ by at most 0.04; the rest of what the gate reads is spread over other directions. We did not predict that the output-space density, once the count is removed, correlates at $+0.42$ and $+0.40$.
\end{itemize}
We conclude that the gate reads a trace of the training interactions in the image expert's output, in part its component along the mean user embedding, not the item's appearance; the rest of that trace is not identified. These are correlations and interventions on two seeds of one augmentation design.

%====================================================================
\section{Limitations, Conclusions, Unresolved Questions}

\textbf{Limitations.} The reproduction is close, not exact (Rec@10, cause not identified). The features are LATTICE's, not CLIP; sizes of the loss and recovery may differ with other features. The noise is simulated on embeddings, never real text. Baby has the authors' split; Sports is our own construction. Three seeds for the early experiments, one seed for sensitivity sweeps; the paired differences are of the order of the seed std and no significance test is claimed. The selected $\gamma$ is the grid minimum, where a large share of VGR's gain is uniform muting of text. The mechanism findings rest on two seeds per setting and one noise family. 

\textbf{Conclusions.} The released code reproduces Rec@20 and NDCG@20 and is 0.0025 short on Rec@10. When 40\% of the cold items' text is corrupted, Recall@20 falls 14\% and the gate keeps trusting the text. VGR recovers about a third of that loss at no measurable clean cost, but about half of its gain comes from muting text, so as a fix it is modest. Blending in the text of the visually nearest warm items recovers much more (about 85\% on five Baby seeds), costs 0.6 s and no measurable clean accuracy, and random-neighbour and permuted-cosine controls gain far less; it holds under shuffle and Gaussian noise, rates 0.2--0.6 and mismatched validation noise. Flagging items with VGR's cosine is a small advantage over the blind blend ($+0.0019$ over ten seeds) that disappears at high noise. On Sports the ranking holds but M3 recovers only about 38\%. The text preference of the augmented gate on cold items follows from items having no training interactions.

\textbf{Unresolved questions.} What else, besides the user-direction component, the gate reads in the image output. Whether augmentation designs other than a lower rate (image-side noise, fixed prototypes) behave differently. A $\gamma$ rule that maximises the validation gain over the shuffled control moves the choice to 2--4 in two of three seeds and raises the cosine-specific gain slightly at the cost of total gain (H24); the differences are within seed noise, so the released rule stays. Whether the pattern holds with real clickbait text, genuine CLIP features and the authors' own splits on other datasets. What closes the Rec@10 gap; the one candidate left is an epoch chosen with the test set in the published run.

\textbf{Errors found and corrected along the way.}
\begin{itemize}
\item \emph{VGR was applied in training} in the first version of the \code{improved} mode; it became an evaluation-time switch before any reported result.
\item \emph{The CLIP assumption} did not fit the shapes of the distributed files; it became a recorded deviation and later a measured one (identical to LATTICE's).
\item \emph{The H3 verdict} was first written ``supported'' although 5.6\% exceeded the 5\% limit; it now reads supported for quality, not latency.
\item \emph{The reproducibility claim} (``a rerun reproduces every number'') failed after the system PyTorch changed (2.14.1 against the pinned 2.9.1); reinstalling 2.9.1 in a separate environment reproduced all saved baselines exactly, so the claim now names the environment.
\item \emph{The oracle} was labelled an upper bound (M3 and M1 beat it); \emph{the H17 text} named negative sampling wrongly; mechanism guesses (bias, norm, distance) were rejected by measurement.
\end{itemize}

%====================================================================
% Bibliography is embedded so the build needs no .bib file.
% Venue, pages and DOI were checked against the Crossref record or the publisher page.
\begin{thebibliography}{9}
\bibitem{nguyen2025mamex} V.-K. Nguyen, D.-H. Pham, H.-S. Nguyen, C.-V. T. Nguyen, H.-Q. Le, and D.-T. Le, ``{MAMEX}: Multi-modal adaptive mixture of experts for cold-start recommendation,'' in \emph{Proc. 34th ACM Int. Conf. Information and Knowledge Management (CIKM)}, 2025, pp. 5053--5057, doi:10.1145/3746252.3760837.
\bibitem{he2016vbpr} R. He and J. McAuley, ``{VBPR}: Visual Bayesian personalized ranking from implicit feedback,'' in \emph{Proc. AAAI Conf. Artificial Intelligence}, vol.~30, no.~1, 2016, pp. 144--150, doi:10.1609/aaai.v30i1.9973.
\bibitem{wei2019mmgcn} Y. Wei, X. Wang, L. Nie, X. He, R. Hong, and T.-S. Chua, ``{MMGCN}: Multi-modal graph convolution network for personalized recommendation of micro-video,'' in \emph{Proc. 27th ACM Int. Conf. Multimedia (MM)}, 2019, pp. 1437--1445, doi:10.1145/3343031.3351034.
\bibitem{wei2021clcrec} Y. Wei, X. Wang, Q. Li, L. Nie, Y. Li, X. Li, and T.-S. Chua, ``Contrastive learning for cold-start recommendation,'' in \emph{Proc. 29th ACM Int. Conf. Multimedia (MM)}, 2021, pp. 5382--5390, doi:10.1145/3474085.3475665.
\bibitem{liu2026motorec} J. Liu, Z. Zhang, and R. C. C. Cheung, ``{MoToRec}: Sparse-regularized multimodal tokenization for cold-start recommendation,'' in \emph{Proc. 40th AAAI Conf. Artificial Intelligence}, vol.~40, no.~18, 2026, pp. 15324--15332, doi:10.1609/aaai.v40i18.38558.
\bibitem{bai2024milk} H. Bai, L. Wu, M. Hou, M. Cai, Z. He, Y. Zhou, R. Hong, and M. Wang, ``Multimodality invariant learning for multimedia-based new item recommendation,'' in \emph{Proc. 47th Int. ACM SIGIR Conf. Research and Development in Information Retrieval}, 2024, doi:10.1145/3626772.3658596.
\bibitem{kim2025dgmrec} J. Kim, H. Kang, S. Kim, K. Kim, and C. Park, ``Disentangling and generating modalities for recommendation in missing modality scenarios,'' in \emph{Proc. 48th Int. ACM SIGIR Conf. Research and Development in Information Retrieval}, 2025, pp. 1820--1829, doi:10.1145/3726302.3729953.
\bibitem{mmsrec} kz-song, ``{MMSRec} code repository.'' [Online]. Available: \url{https://github.com/kz-song/MMSRec} (README read 2026-10-23).
\bibitem{radford2021clip} A. Radford \emph{et al.}, ``Learning transferable visual models from natural language supervision,'' in \emph{Proc. 38th Int. Conf. Machine Learning (ICML)}, PMLR vol.~139, 2021, pp. 8748--8763.
\bibitem{zhang2021lattice} J. Zhang, Y. Zhu, Q. Liu, S. Wu, S. Wang, and L. Wang, ``Mining latent structures for multimedia recommendation,'' in \emph{Proc. 29th ACM Int. Conf. Multimedia (MM)}, 2021, pp. 3872--3880, doi:10.1145/3474085.3475259.
\end{thebibliography}

%====================================================================
\appendices

\section{AI-Use Disclosure}
\label{app:ai}

Tool: Claude Code (Anthropic CLI), model Claude Sonnet 5.5, used for this project from 2026-10-16, including web search and page fetching. The git commit carries no AI trailer, so this appendix is the record of AI contribution. The author remains responsible for reading the sources, checking the citations and defending the conclusions.

\begin{table}[!h]
\centering\footnotesize
\begin{tabularx}{\columnwidth}{@{}p{1.8cm}X@{}}
\toprule
Area & Contribution \\
\midrule
Code & Restructured the upstream code into a package; wrote the noise simulator, VGR and its controls, $\gamma$ selection, evaluation report, CLI modes, plotting, table and failure-analysis scripts, the experiment scripts and the tests (test-first); translated Vietnamese and Chinese comments to English \\
Design & Proposed the noise protocol, the controls and the $\gamma$-selection rule; wrote the hypotheses before full-training results were read \\
Literature & Ran the searches, read paper pages, compiled the matrix and bibliography; checked DOIs, venues and pages against Crossref \\
Debugging & Found that the \code{improved} mode also trained with VGR and changed it \\
Runs & Launched and monitored the experiment grid and the diagnostic runs; generated tables, figures and the failure analysis from saved outputs; checked that a repeated run reproduces them \\
Report & Drafted the report, README and reproduction notes; every number comes from the saved experiment outputs or the cited papers; compiled the PDF locally and fixed layout warnings \\
Review & Checked the report against the course brief; compared the CLIP embedding dimensions with the CLIP paper's appendix tables (the author ran the \code{open_clip} listing) \\
\bottomrule
\end{tabularx}
\end{table}

\section{Literature Comparison Matrix}
\label{app:matrix}

One row per paper, one column per field of the course template. [Authors]: reported by the paper; [Interpretation]: our inference. MAMEX is the paper reproduced.

\clearpage
\onecolumn
\begingroup\centering\scriptsize\setlength{\tabcolsep}{4pt}
\renewcommand{\arraystretch}{1.1}
\begin{longtable}{@{}p{1.7cm}*{7}{>{\raggedright\arraybackslash}p{2.0cm}}@{}}
\caption{Literature comparison matrix.}\label{tab:matrix}\\
\toprule
Paper & Question \& contribution & Modalities \& assumptions & Mechanism \& objective & Data \& protocol & Baselines \& evidence & Resources \& reproducibility & Limitations \& connections \\
\midrule
\endfirsthead
\toprule
Paper & Question \& contribution & Modalities \& assumptions & Mechanism \& objective & Data \& protocol & Baselines \& evidence & Resources \& reproducibility & Limitations \& connections \\
\midrule
\endhead
\bottomrule
\endfoot
VBPR~\cite{he2016vbpr} & Can product-image appearance improve implicit-feedback ranking, including for items with little feedback? [Authors]: adds visual factors to BPR-MF. & Image (4096-d CNN FC7, ImageNet-pretrained) and implicit feedback. Items need only an image to be scored. Users with $<5$ positives dropped. & $\hat{x} = \alpha + \beta_u + \beta_i + \gamma_u^{\top}\gamma_i + \theta_u^{\top}(E f_i) + \beta'^{\top} f_i$; BPR-OPT; 20 factors, half visual. & Amazon Women, Men, Phones, Tradesy. One random item per user for val and test. Cold = test items with $<5$ training positives (not zero). AUC. & RAND, MP, IBR, MM-MF, BPR-MF; hyperparameters tuned on validation; no seeds or significance tests. Women cold-start AUC: VBPR 0.6813 vs BPR-MF 0.5281; Phones: VBPR 0.6056 below IBR 0.6319 [Authors]. & Desktop, 4 cores, 32 GB; about 3.5 h on Women. Code ``at publication time'', no link in the text. & Single image per item; visual/latent split fixed 50/50; no component ablation. Image only. \\
\midrule
MMGCN~\cite{wei2019mmgcn} & Model modality-specific user preferences for micro-video recommendation. [Authors]: one user-item graph per modality with message passing. & Visual, acoustic, textual pre-extracted features. Items must already sit in the interaction graph, so it is not a cold-item method [Interpretation]. & Per-modality GCN aggregation (mean or max) with combination layers and a shared user-ID embedding; representations are the sum of the modality-specific ones, scored by inner product; BPR loss with L2; SGD. & Tiktok 726,065 / Kwai 1,664,305 / MovieLens 1,239,508 interactions. Random 8:1:1 over interactions; each test pair ranked against 1000 unobserved videos; Precision, Recall, NDCG@10. & VBPR, ACF, GraphSAGE, NGCF. Recall gains over the strongest baseline 15.59\% / 10.23\% / 8.76\% (Kwai / Tiktok / MovieLens) [Authors]. Multimodal beats each single modality; visual is most effective, text least on Kwai and Tiktok. & Authors state they release codes, parameters and baselines. & Random interaction split: no cold-item evaluation, no corrupted-modality test. Calls the Kwai/Tiktok descriptions ``noisy, incomplete, and even irrelevant'' but runs no corrupted-text test. \\
\midrule
CLCRec~\cite{wei2021clcrec} & Cold-start item embeddings from an information-theoretic view: maximize mutual information between content and collaborative signals. & Pre-extracted visual, acoustic, text features (ResNet, VGGish, Sentence2Vector); Amazon uses 64-d visual only. Cold items scored from the content encoder alone. & Contrastive U-I and R-E objectives, $\lambda L_{RE} + (1-\lambda) L_{UI}$; hybrid training swaps in content features with probability $\rho$. & Movielens, Tiktok, Kwai, Amazon; cold items split 1:1 into val/test, warm items 8:1:1. Recall@10, NDCG@10. & DUIF, DropoutNet, MTPR, CB2CF, Heater; same 64-d embedding for all. Single runs, no variance. & Code and data released; hardware not stated. & Single runs without variance. No text-corruption test. \\
\midrule
MoToRec~\cite{liu2026motorec} & Argues that aligning noisy continuous multimodal embeddings causes ``semantic fog''; proposes discrete semantic tokens. & Visual (BEiT), text (BGE), learnable ID embeddings. Cold = items with $<10$ training interactions. Deployment for new items not described [Authors' text]. & Sparsity-regularised RQ-VAE per modality, rarity weighting, LightGCN-based multi-source graph encoder; $L_{BPR}+L_{CL}+L_{RQ}+L_{reg}$. & Amazon Baby (19,445 users, 7,050 items), Sports, Clothing. Random 8:1:1 over interactions; cold group = test items with $<10$ training interactions. R, N@\{10,20\}. & MF-BPR, LightGCN, VBPR, MMGCN, LATTICE, FREEDOM, BM3, LGMRec and others, in MMRec; grid search only described for MoToRec [Interpretation: baseline tuning fairness unverifiable]. Baby R@20 0.1077; Baby cold-start R@20 0.0347 [Authors]. & Hardware and code not stated. & Cold-start differs from complete cold start (items still have some interactions). Claims noisy embeddings are a problem but does not corrupt inputs to test it. \\
\midrule
MAMEX~\cite{nguyen2025mamex} & Cold-start fusion fails because concatenation, averaging or fixed weights ignore modality relations; adaptive gating across modality experts. & CLIP image and text features [Authors]; no dimensions stated. Complete item cold start: all interactions of dev/test items removed. & Modality-specific MoE (linear experts, top-$k$) and sparse-softmax gate over $[z_{img}\Vert z_{text}]$; $L_{BPR}+\lambda_1 L_{align}+\lambda_2 L_{adapter}+\lambda_3 L_{fusion}+\lambda_4\lVert\Theta\rVert^2$. & Amazon Baby, Clothing, Sports; 8:1:1; Rec, NDCG@\{10,20\}. & MTPR, AlignRec, MetaEmbed, DropoutNet, CLCRec, GAR, MILK. Baby: Rec@20 0.0771, NDCG@20 0.0363 [Authors]. Modality ablation (figure only): text beats vision, both beat either alone. & Code public (\url{github.com/L2R-UET/MAMEX}); batch size, epochs, embedding size and hardware not stated. & No limitations section; future work lists missing modalities. Clean features only; the gate's behaviour under corrupted text is untested. Direct target of this project. \\
\midrule
DGMRec~\cite{kim2025dgmrec} & Multimodal recommenders degrade when item modalities are missing, and filling in a global or nearest-neighbour mean does not prevent it [Authors]. Disentangles general and modality-specific features and generates the missing ones. & Image and text on Amazon (a CNN for images, SBERT for text [Authors]); TikTok adds audio. Some or all modality features of an item are absent; present features are assumed correct [Interpretation]. & Autoencoder; general and specific encoders trained with two information-based losses; a missing modality is generated from the other modalities and from the modality preferences of the item's users; generated features refine the item-item graph. & Amazon Baby, Sports, Clothing and TikTok; 8:1:1; Recall, NDCG@\{20,50\}. Missing setting: a third of items miss 0, 1 or 2 modalities; new-item setting: 20\% of items only in test; missing ratios 0 to 80\%. & 5 CF models, 10 multimodal recommenders (VBPR, MMGCN, LATTICE, LGMRec, GUME and others) and 3 missing-modality-aware ones; best in the missing and new-item settings; drops 10.2\% vs 15.1\% for the compared baseline at an 80\% missing ratio [Authors]. & Code public (\url{github.com/ptkjw1997/DGMRec}); Adam, embedding size 64, up to 1,000 epochs, early stopping 30, $\lambda_1,\lambda_2$ tuned in $\{1,\allowbreak 0.1,\allowbreak 0.01,\allowbreak 0.001\}$. & No limitations section. Handles absent modalities, not present but misleading ones; the injected features it criticises are shown to hurt, the closest it comes to a corrupted input, but no corrupted-text test is run [Interpretation]. \\
\midrule
MILK~\cite{bai2024milk} & New-item recommendation when items miss a modality. [Authors]: invariant learning, so that the user's preference prediction does not change across modality-missing environments; cross-modality alignment plus cyclic mixup. & Visual and text on Baby and Clothing; TikTok adds audio. Pre-extracted features, extractor and dimensions not stated in the text read [Interpretation]. A missing modality is mean-filled; present features are assumed correct [Interpretation]. & Modality-specific extractors with alignment functions between modalities; fusion weights $\Theta_1\sim$ Dirichlet, $M$ cyclic shifts and an equal-weight environment; $L = E_e L_e + \beta\,\mathrm{Var}_e L_e + \lambda L_{align} + \lVert\Phi\rVert^2$, BPR inside each environment. & Amazon Baby (19,442 users; 5,640 train, 705 val, 705 test items), Clothing, TikTok. 20\% of items new. FTMT: full training, one modality missing at test; MTMT: 30\% of training items also miss one. Recall, NDCG@\{10,20\}. & DUIF, MICRO, DropoutNet, MTPR, Heater, CLCRec, CCFCRec, GAR, GoRec. Baby FTMT Rec@20 0.0628, NDCG@20 0.0277 (best baseline GoRec 0.0544, 0.0249); mean of five runs, no variance stated in the text read [Authors]. MAMEX lists MILK at Rec@20 0.0730 on Baby, a different setting [Interpretation]. & PyTorch, 64-d, batch 2048, Adam 0.001, early stopping; code public (\url{github.com/HaoyueBai98/MILK}); hardware not stated. $\beta$ is tuned per dataset (1000 on Baby). & Absence, not misleading content, is tested. Same item split as MAMEX and as this project (705 val and 705 test items). MAMEX's README acknowledges MILK for foundational concepts. \\
\end{longtable}
\endgroup
\twocolumn

\section{Reproduction Record}
\label{app:record}

\textbf{Source.} Reproduced code: \url{https://github.com/L2R-UET/MAMEX}, commit \code{88a932eca821a5657b2b7fc0b6445e4685f4dfab}. This repository (URL on page 1) restructures it; model, loss and training loop are unchanged except for two switches that default to off. Commit timestamps do not show when runs were made (the history was built after the experiments); the tables regenerate identically from the saved outputs.

\textbf{Environment.} Linux 7.0.0; Python 3.13.9; PyTorch 2.9.1 (CUDA 12.8) pinned in a separate virtual environment, which reproduces the saved baselines of seeds 2025--2027 exactly (the system PyTorch 2.14.1+cu130 does not); Intel Core i9-12900K, 62 GB RAM, NVIDIA RTX A4000. Packages are pinned in the requirements file.

\textbf{Data.} Baby, item cold-start split as distributed by the authors (\url{https://drive.google.com/file/d/1C-18Y84lMS5xsRGwKQYa9P0Tot-gxgYY/view}; no raw text, images or ASIN mapping). ASINs of the 7,050 items come from the \code{item\_list.txt} of the LATTICE Baby files; the demo's item browser joins them to the 2023 Amazon catalogue (titles and image links for 4,353 items). Data are not redistributed; Amazon review data are for research under their own terms.

\begin{table}[!h]
\centering\tiny\setlength{\tabcolsep}{1pt}
\caption{SHA-256 of the data files used (7,050 items; image 4096-d float64, text 1024-d float32).}
\begin{tabular}{@{}ll@{}}
\toprule
File & SHA-256 \\
\midrule
train.txt & \texttt{330c8bdff9b82a3b9b83f24dae5830c915a9aff239352f788959a15b71c3b83f} \\
val.txt & \texttt{4768ff30900b4ff84d053ee2df416cb9b51248a6b20eed1fc139f7f47ed82382} \\
test.txt & \texttt{d72c5f69bbe1a3ae6683efecdb6ed40e24613c05dc0d5f1617f40dc79be18aea} \\
cold\_item\_index.npy & \texttt{fd57f2947047eed66be739be65355b0214c129d0d6e31fa55ecf95e3d3b20abd} \\
image\_feat.npy & \texttt{2fc1695bc4fdf58a13d325dd0a63fe270d864ec4c8b26b780f8a1eee9682747e} \\
text\_feat.npy & \texttt{c0c9b38fc9871d635814484e5948f4c4eeb2e0c33668babdf3747898ce4400b2} \\
\bottomrule
\end{tabular}
\end{table}

\textbf{Seeds and configuration.} Training seeds 2025--2027 (first experiments, Sports), 2025--2034 (ten-seed comparison). The corruption seed defaults to the training seed; the validation mask uses the seed plus 10,000. Upstream defaults otherwise (learning rate $10^{-3}$, batch 2,048, 128-d, early stop 200).

\textbf{Commands} (repository root; full list in \code{docs/reproducibility.md}):
{\raggedright\scriptsize\ttfamily
PYTHONPATH=src python -m pytest -q\ \ (158 tests, no GPU)\\
./scripts/demo.sh\ \ (about 2 min)\\
DATA\_PATH=processed\_data\_baby/cold\\
\ \ SEEDS="2025 2026 2027" ./scripts/run\_experiments.sh\\
./scripts/run\_repair.sh\ \ (repairs, H8--H16)\\
SRC=<Sports files> ./scripts/run\_second\_dataset.sh\par}
Tables, failure analysis, ranking metrics and figures come from the reporting modules run on the saved outputs.

\textbf{Cost.} About 19 min of training per run alone, on one GPU. The first 54 result files (core grid, sweeps, 16 diagnostic runs, ablation, gate and representation probes) took about 42 hours of training with several jobs sharing the GPU. The checkpoint, curve and later runs (H11 to H18: seeds 2028--2034 with checkpoints, per-epoch curves, augmented checkpoints, held-out runs, Sports) add 40.1 hours. The repair evaluations, interventions and analyses need no training. Checkpoints are not published; every result comes from retraining with the stated seeds.

\end{document}
